\section{Related Work}

Our work extends code generation benchmarks with process metrics measuring debugging efficiency. We position our framework relative to three areas: correctness evaluation, execution feedback for training, and multi-test benchmarks.

\subsection{Code Generation Benchmarks}

HumanEval~\cite{Chen2021Evaluating} and MBPP~\cite{Austin2021Program} established pass@k as the standard correctness metric for code generation, measuring the probability that at least one of k samples passes unit tests. These benchmarks evaluate generation quality but not iterative debugging — problems provide 1-5 test cases, insufficient to reveal strategic behavior. Extensions like HumanEval+~\cite{Liu2023Is} improve test coverage to detect brittleness, but retain outcome-focused evaluation (did the code pass?) without analyzing debugging trajectories.

Competitive programming benchmarks (CodeContests~\cite{Li2022Competition}, APPS~\cite{Hendrycks2021Measuring}) use multi-test problems (10-100 test cases) where strategic debugging becomes measurable. However, existing work reports final solve rates or test-pass percentages~\cite{Jain2024LiveCodeBench}, not how efficiently agents reach these outcomes. Our framework analyzes debugging trajectories on such multi-test problems, measuring fix-impact-ratio (tests passed per modification) and error clustering — metrics invisible to total pass rate.

\subsection{Execution Feedback in Code Generation}

Execution-based feedback improves code generation through reinforcement learning~\cite{Le2022CodeRL} or iterative refinement with error messages~\cite{Chen2023Teaching}. These approaches use feedback for \emph{training} (e.g., RL reward signals from test execution) rather than \emph{evaluation} — they measure whether feedback improves generation, not how agents strategically utilize feedback during debugging.

Our work is orthogonal: we measure strategic debugging behavior (clustering errors, prioritizing fixes) as an evaluation framework, applicable to any agent regardless of training method. An RL-trained agent and a prompted GPT-4 agent can both be evaluated with fix-impact-ratio to quantify debugging efficiency — the metric measures process, not training paradigm.

\subsection{Multi-Test Programming Evaluation}

LeetCode and Codeforces provide multi-test problems (15-50 test cases) used in agent evaluation~\cite{Jain2024LiveCodeBench}, but existing work focuses on final correctness rates. SWE-bench~\cite{Jimenez2024SWE} evaluates GitHub issue resolution with repository-level context, measuring patch success without analyzing debugging efficiency when test suites are available.

Our framework could extend to SWE-bench scenarios where test suites exist — measuring how efficiently agents pass tests reveals strategic capability beyond issue resolution rates. Similarly, multi-test competitive programming benchmarks support fix-impact-ratio measurement, but prior work hasn't analyzed debugging trajectories systematically.

\subsection{Positioning}

We extend outcome-focused benchmarks (HumanEval pass@k, LeetCode solve rates) by adding process metrics that reveal \emph{how} agents debug. Our metrics are execution-based (no subjective model judges required), applicable to any multi-test programming scenario (competitive programming, GitHub issues with tests), and orthogonal to training methods (measure behavior, not how the agent was trained). This positions our work as a complementary evaluation dimension: correctness (pass@k) + efficiency (fix-impact-ratio) together provide a fuller picture of agent capability.

The key distinction is that we measure debugging as a \emph{process} with testable stages (clustering $\rightarrow$ prioritization), validated via execution-based metrics (clustering coefficient, fix-impact-ratio) rather than relying on final pass rates alone. This reveals strategic debugging operates as a feedback loop effective within revealed test execution, but not a learning system enabling pattern transfer to unseen tests — a boundary existing benchmarks cannot detect.
